\documentclass[11pt]{article}
\usepackage[margin=26mm]{geometry}
\usepackage{amsmath,amssymb,amsthm,mathtools,microtype,hyperref}
\hypersetup{hidelinks}
\newtheorem{theorem}{Theorem}
\newtheorem{proposition}[theorem]{Proposition}
\newcommand{\ddc}{dd^c}
\newcommand{\Ric}{\operatorname{Ric}}
\newcommand{\vol}{\operatorname{vol}}
\newcommand{\tr}{\operatorname{tr}}
\title{Divergence, rigidity and conformal gauges\newline for Chern--Ricci-flat LCK metrics}
\author{}
\date{6 October 2026}
\begin{document}
\maketitle
\begin{center}\small Scoped metric and conformal-class results; the general existence question remains open.\end{center}
\begin{abstract}
For a compact Chern--Ricci-flat locally conformally K\"ahler metric,
the squared $L^2$ norm of the anti-$J$-invariant covariant derivative
of its Lee form equals half the squared norm of its divergence.
This gives equivalent Gauduchon, holomorphic-Lee and Vaisman conditions
for that metric. An explicit fixed Kodaira surface has a strictly LCK
flat conformal gauge which is not Gauduchon, and its conformal class
contains no metric satisfying both conditions. The same complex surface
has an explicit Vaisman metric. These are metric and conformal-class
results; the unrestricted existence question for an arbitrary compact
ordinary Bott--Chern Calabi--Yau LCK manifold remains unresolved.
\end{abstract}

\section{Conventions and scope}
Let $(X,J,g,\omega)$ be connected and compact, of complex dimension
$n\geq2$, with $\omega(X,Y)=g(JX,Y)$. We use
\[
 d^c=\frac{i}{4\pi}(\bar\partial-\partial),
 \qquad \ddc=\frac{i}{2\pi}\partial\bar\partial,
 \qquad J\alpha=-\alpha\circ J
\]
for real one-forms. Thus $dJdf=4\pi\ddc f$. A strictly LCK metric has
$d\omega=\vartheta\wedge\omega$ with $d\vartheta=0$ and nonexact
$\vartheta$. It is Vaisman when $\nabla\vartheta=0$. The Chern--Ricci
form is $\rho(g)=-i\partial\bar\partial\log\det(g)$; its normalized
representative $\rho(g)/(2\pi)$ represents ordinary $c_1^{\rm BC}(T_X)$.
All divergences, norms and integrals below use this same $g$.

The same-metric Gauduchon and Chern--Ricci-flat rigidity theorem is
established in \cite[Theorem A]{BO}. We retain the divergence terms
in its connection, Bianchi and Weitzenb\"ock argument. Ordinary
Bott--Chern vanishing alone does not supply a simultaneous gauge.
No theorem for an already Vaisman manifold is applied to a general
LCK manifold, and no deformation or cover replaces the given complex
structure.

\section{The integral identity and its rigidity consequences}
Set
\[
 A=\vartheta^\sharp,\quad B=JA,\quad T=\nabla\vartheta,\quad
 P(X,Y)=\tfrac12\{T(X,Y)+T(JX,JY)\},\quad T^-=T-P,
\]
\[
 q=d_g^*\vartheta=-\tr_gT,\qquad u=|\vartheta|_g^2.
\]
Closedness of $\vartheta$ makes $T$ symmetric; $P$ and $T^-$ are
orthogonal, so $|T|^2=|P|^2+|T^-|^2$.

\begin{theorem}\label{th:rigidity}
For every compact Chern--Ricci-flat strictly LCK metric,
\begin{equation}\label{eq:identity}
 \int_X|T^-|^2\,d\vol_g=\frac12\int_X|q|^2\,d\vol_g.
\end{equation}
For this fixed metric the following conditions are equivalent:
Gauduchon; $q=0$; $A$ is real holomorphic; and Vaisman.
The equality \eqref{eq:identity} is an integral identity, rather than
an asserted pointwise identity for general metrics.
\end{theorem}

\begin{proof}
Locally $\vartheta=df$ and $\widetilde g=e^{-f}g$ is K\"ahler. Its
Levi--Civita connection is
\[
 D_XY=\nabla_XY-\tfrac12\vartheta(X)Y
                 -\tfrac12\vartheta(Y)X+\tfrac12g(X,Y)A.
\]
The equation $DJ=0$ gives
\begin{equation}\label{eq:nablaJ}
 (\nabla_XJ)Y=\tfrac12\{\vartheta(JY)X-\vartheta(Y)JX
                  -g(X,JY)A+g(X,Y)JA\}.
\end{equation}
Differentiating $J\vartheta=-\vartheta\circ J$ and using symmetry of
$T$ consequently gives
\[
 dJ\vartheta(X,JY)=2P(X,Y)-u g(X,Y)
              +\vartheta(X)\vartheta(Y)+J\vartheta(X)J\vartheta(Y).
\]
Flatness and the local conformal volume change give
$\rho(\widetilde g)=n i\partial\bar\partial f=(n/2)dJ\vartheta$.
Contracting the curvature difference of $D$ and $\nabla$ gives the
real conformal Ricci formula
\[
 \Ric_D=\Ric_g+(n-1)T+\tfrac{n-1}{2}\vartheta\otimes\vartheta
                       +\tfrac12\{-q-(n-1)u\}g.
\]
The K\"ahler Ricci form of $\widetilde g$ equals its real Ricci tensor
paired with $J$. Combining these expressions yields
\begin{equation}\label{eq:Ric}
 R:=\Ric_g=nP-(n-1)T+\tfrac{q-u}{2}g
               +\tfrac12\vartheta\otimes\vartheta
               +\tfrac n2J\vartheta\otimes J\vartheta.
\end{equation}
Its trace is
$s=(n-1)q-(n-1)u/2$. The local formulas agree with
\cite[Section 4, equations (2)--(4)]{BO}; no $q=0$ hypothesis has
been inserted into \eqref{eq:Ric}.

Write $\|\cdot\|$ for the $L^2$ norm and put
$E_A=\int_XT(A,A)d\vol_g$, $E_B=\int_XT(B,B)d\vol_g$.
Since $T(A,A)=A(u)/2$ and $\operatorname{div}A=-q$,
\begin{equation}\label{eq:uq}
 \int_Xuq\,d\vol_g=2E_A.
\end{equation}
Pair the contracted Bianchi identity
$\delta R+ds/2=0$ with $\vartheta$ and integrate. Here $\delta$ is
minus covariant contraction, adjoint to symmetric covariant derivative.
It gives $\int\langle R,T\rangle+\frac12\int s q=0$.
Substitution of \eqref{eq:Ric} and \eqref{eq:uq} gives
\begin{equation}\label{eq:Bianchi}
 0=n\|P\|^2-(n-1)\|T\|^2+\tfrac{n-2}{2}\|q\|^2
                    +\tfrac{4-n}{2}E_A+\tfrac n2E_B.
\end{equation}
The integrated Weitzenb\"ock identity for a closed one-form gives
\[
 \|q\|^2=\|T\|^2+\int_XR(A,A)\,d\vol_g.
\]
Equation \eqref{eq:Ric} has
$R(A,A)=(1-n/2)T(A,A)+(n/2)T(B,B)+qu/2$; hence
\begin{equation}\label{eq:Bochner}
 \|q\|^2=\|T\|^2+\tfrac{4-n}{2}E_A+\tfrac n2E_B.
\end{equation}
Eliminating $E_A,E_B$ between \eqref{eq:Bianchi} and
\eqref{eq:Bochner} proves \eqref{eq:identity}.

For the Gauduchon equivalence, the LCK equation and the Hodge identity
$*\vartheta=J\vartheta\wedge\omega^{n-1}/(n-1)!$ give
\[
 d^c\omega^{n-1}=\frac{n-1}{4\pi}
                       J\vartheta\wedge\omega^{n-1},\qquad
 \ddc\omega^{n-1}=-\frac{n-1}{4\pi n}q\,\omega^n.
\]
Thus Gauduchon is exactly $q=0$. In this case \eqref{eq:identity}
forces $T=P$, so $T(B,B)=T(A,A)$. Equation \eqref{eq:uq} gives
$E_A=E_B=0$, and \eqref{eq:Bochner} forces $T=0$.

Equation \eqref{eq:nablaJ} gives $\nabla_AJ=0$. Let $S$ be the
symmetric endomorphism with $g(SX,Y)=T(X,Y)$. Then
\[
 (\mathcal L_AJ)Y=JSY-SJY.
\]
Consequently $A$ is real holomorphic exactly when $T=P$.
Equation \eqref{eq:identity} then gives $q=0$, hence $T=0$.
Conversely $T=0$ gives both other conditions. Under strictness the
parallel Lee form is nonzero, as required for Vaisman.
\end{proof}

\section{An explicit compact example with incompatible gauges}
\begin{proposition}\label{prop:example}
On one fixed compact complex Kodaira surface there is a strictly LCK
Chern--Ricci-flat metric which is not Gauduchon. Its conformal class
contains no simultaneously Gauduchon and Chern--Ricci-flat metric.
The same surface has an explicit Vaisman metric, so this is not a
counterexample to the unrestricted P05 existence question.
\end{proposition}

\begin{proof}
On $G=\mathbb R^4$ use the multiplication
\[
 (x,y,z,t)(x',y',z',t')=(x+x',y+y',z+z'+xy',t+t').
\]
The integer subgroup $\Gamma=\mathbb Z^4$ is a lattice: its left
action is $(x,y,z,t)\mapsto(x+a,y+b,z+c+ay,t+d)$.
Reduction of $x,y,t$ and then $z$ to unit intervals supplies a compact
fundamental set. Thus $X=\Gamma\backslash G$ is connected and compact.
The global forms
\[
 e^1=dx,\quad e^2=dy,\quad e^3=dz-xdy,\quad e^4=dt
\]
have $de^3=-e^{12}$ and the other derivatives zero. Declare
$\phi^1=e^1+i e^2$, $\phi^2=e^3+i e^4$ to be of type $(1,0)$.
Then $d\phi^1=0$ and
$d\phi^2=-(i/2)\phi^1\wedge\overline{\phi^1}$ have no $(0,2)$
part, so the complex structure is integrable.
$\Sigma=\phi^1\wedge\phi^2$ is closed and nowhere vanishing,
hence holomorphic. The canonical bundle is holomorphically trivial,
which proves ordinary $c_1^{\rm BC}(T_X)=0$.

Take a unit-periodic $f(x)$ with $C=1-f''>0$, and set
\[
 \vartheta_1=-e^4+f'e^1,\quad \eta_1=J\vartheta_1=e^3+f'e^2,
 \quad \omega_1=\vartheta_1\wedge\eta_1+C e^{12}.
\]
This is the positive Hermitian form of
$g_1=\vartheta_1^2+\eta_1^2+C\{(e^1)^2+(e^2)^2\}$.
One has $d\vartheta_1=0$, $d\eta_1=-C e^{12}$ and
$d\omega_1=\vartheta_1\wedge\omega_1$.
The $t$-circle has Lee period $-1$, so the metric is strict.
Its Lee dual is $-\partial_t$. All coefficients are independent of
$t$, making this field Killing; closedness of $\vartheta_1$ makes
$\nabla\vartheta_1$ symmetric and Killing makes it skew. Thus $g_1$
is Vaisman on this fixed complex manifold.

Its volume is $\omega_1^2/2=C e^{1234}$, while
$\Sigma\wedge\overline\Sigma$ has constant density in $e^{1234}$.
Put $h=C^{-1/2}$ and $\omega=h\omega_1$. Then
$\omega^2/2=e^{1234}$, so its Chern--Ricci form is zero. Its Lee form
is $\vartheta=-e^4+d(f+\log h)$, retaining period $-1$.
If $E_i$ are dual to $e^i$, its Lee dual is
\[
 A=h'E_1+(f'h'-h^{-1})E_4.
\]
The metric volume is $dx\,dy\,dz\,dt$. Each $E_i$ has zero
divergence in this volume, hence $\operatorname{div}A=h''$ and
$q=d_g^*\vartheta=-h''$.

For the explicit choice
\[
 f(x)=\frac{\cos(2\pi x)}{16\pi^2},\qquad
 C(x)=1+\tfrac14\cos(2\pi x)\geq\tfrac34,
\]
direct differentiation gives
\[
 q(0)=-\frac{4\sqrt5\pi^2}{25},\qquad
 q(1/2)=\frac{4\sqrt3\pi^2}{9}.
\]
Thus $\omega$ is not Gauduchon, and Theorem~\ref{th:rigidity}
also shows it is not Vaisman.

For any global real $v$,
$\rho(e^v\omega)=\rho(\omega)-n i\partial\bar\partial v$.
Two flat metrics in a single conformal class therefore have a real
pluriharmonic logarithmic ratio, constant on a connected compact
complex manifold by the maximum principle. The only flat metrics
conformal to $\omega$ are its positive constant multiples; constant
scaling cannot make its nonzero $q$ vanish. This proves the asserted
obstruction. The already exhibited $\omega_1$ is a Vaisman and
Gauduchon representative in that class, with generally nonzero
Chern--Ricci form.
\end{proof}

\section{The remaining existence question}
If ordinary $c_1^{\rm BC}(T_X)=0$, write
$\rho(\omega)/(2\pi)=\ddc F$ for a global smooth $F$.
Then $e^{F/n}\omega$ is Chern--Ricci flat and its Lee form is
$\vartheta+dF/n$, still nonexact. This proves existence and
uniqueness up to scale of the flat conformal gauge, rather than
existence of a metric with $q=0$ in it. The explicit example
shows why these assertions cannot be combined with the separate
Gauduchon conformal theorem to supply the same-metric rigidity
hypothesis. It supplies no obstruction to all Vaisman metrics on
$X$, and the arbitrary same-complex-manifold existence problem
remains open in this argument.
\begin{thebibliography}{9}
\bibitem{BO} G. Barbaro and A. Otiman,
\emph{Calabi--Yau locally conformally K\"ahler manifolds},
\href{https://arxiv.org/abs/2509.18364v2}{arXiv:2509.18364v2},
Theorem A and Section 4.
\end{thebibliography}
\end{document}
